
\documentclass{article}
\usepackage{colortbl}
\usepackage{makecell}
\usepackage{multirow}
\usepackage{supertabular}

\begin{document}

\newcounter{utterance}

\centering \large Interaction Transcript for game `matchit\_ascii', experiment `same\_grid', episode 2 with mistral{-}large{-}3{-}azure.
\vspace{24pt}

{ \footnotesize  \setcounter{utterance}{1}
\setlength{\tabcolsep}{0pt}
\begin{supertabular}{c@{$\;$}|p{.15\linewidth}@{}p{.15\linewidth}p{.15\linewidth}p{.15\linewidth}p{.15\linewidth}p{.15\linewidth}}
    \# & $\;$A & \multicolumn{4}{c}{Game Master} & $\;\:$B\\
    \hline

    \theutterance \stepcounter{utterance}  
    & & \multicolumn{4}{p{0.6\linewidth}}{
        \cellcolor[rgb]{0.9,0.9,0.9}{
            \makecell[{{p{\linewidth}}}]{
                \texttt{\tiny{[P1$\langle$GM]}}
                \texttt{Está a participar num jogo de adivinhação colaborativo.} \\
\texttt{O objetivo é descobrir se esta grelha e outra imagem que só eu consigo ver são iguais.} \\
\texttt{Por favor, descreva primeiro a sua grelha. Depois, eu fornecerei a minha descrição e podemos fazer perguntas um ao outro sobre as grelhas para descobrir se são iguais. Agora comece a sua breve descrição da grelha com "DESCRIÇÃO:" seguido da descrição. Não acrescente mais nada. Esta é a sua grelha:} \\
\texttt{X ▢ ▢ ▢ X} \\
\texttt{X X ▢ X X} \\
\texttt{X ▢ X ▢ X} \\
\texttt{X ▢ ▢ ▢ X} \\
\texttt{X ▢ ▢ ▢ X} \\
            }
        }
    }
    & \\ \\

    \theutterance \stepcounter{utterance}  
    & \multicolumn{4}{p{0.6\linewidth}}{
        \cellcolor[rgb]{0.9,0.9,0.9}{
            \makecell[{{p{\linewidth}}}]{
                \texttt{\tiny{[P1$\rangle$GM]}}
                \texttt{DESCRIÇÃO:} \\
\texttt{A grelha é uma matriz 5x5 composta por "X" e quadrados vazios (▢). A estrutura forma um padrão simétrico:} \\
\texttt{{-} **1ª linha**: X ▢ ▢ ▢ X} \\
\texttt{{-} **2ª linha**: X X ▢ X X} \\
\texttt{{-} **3ª linha**: X ▢ X ▢ X} \\
\texttt{{-} **4ª linha**: X ▢ ▢ ▢ X} \\
\texttt{{-} **5ª linha**: X ▢ ▢ ▢ X} \\
\\ 
\texttt{As bordas laterais (1ª e 5ª colunas) são todas "X", exceto o centro da 3ª linha (posição 3,3). Há um "X" adicional no centro da 2ª linha (posição 2,2 e 2,4).} \\
            }
        }
    }
    & & \\ \\

    \theutterance \stepcounter{utterance}  
    & & & \multicolumn{2}{p{0.3\linewidth}}{
        \cellcolor[rgb]{0.9,0.9,0.9}{
            \makecell[{{p{\linewidth}}}]{
                \texttt{\tiny{[GM$|$GM]}}
                \texttt{continue} \\
            }
        }
    }
    & & \\ \\

    \theutterance \stepcounter{utterance}  
    & & \multicolumn{4}{p{0.6\linewidth}}{
        \cellcolor[rgb]{0.9,0.9,0.9}{
            \makecell[{{p{\linewidth}}}]{
                \texttt{\tiny{[GM$\rangle$P2]}}
                \texttt{Está a participar num jogo de adivinhação colaborativo.} \\
\texttt{O objetivo é descobrir se esta grelha e outra imagem que só eu consigo ver são iguais.} \\
\texttt{Por favor, descreva primeiro a sua grelha. Depois, eu fornecerei a minha descrição e podemos fazer perguntas um ao outro sobre as grelhas para descobrir se são iguais. Agora comece a sua breve descrição da grelha com "DESCRIÇÃO:" seguido da descrição. Não acrescente mais nada. Esta é a sua grelha:} \\
\texttt{X ▢ ▢ ▢ X} \\
\texttt{X X ▢ X X} \\
\texttt{X ▢ X ▢ X} \\
\texttt{X ▢ ▢ ▢ X} \\
\texttt{X ▢ ▢ ▢ X} \\
            }
        }
    }
    & \\ \\

    \theutterance \stepcounter{utterance}  
    & & & \multicolumn{4}{p{0.6\linewidth}}{
        \cellcolor[rgb]{0.9,0.9,0.9}{
            \makecell[{{p{\linewidth}}}]{
                \texttt{\tiny{[GM$\langle$P2]}}
                \texttt{DESCRIÇÃO: A grelha é uma matriz 5x5 composta por "X" e quadrados vazios (▢). A estrutura forma um padrão simétrico com "X" nas bordas externas (primeira e última linha, primeira e última coluna), exceto nas posições centrais das linhas 2 a 5. Na segunda linha, há "X" nas colunas 1, 2, 4 e 5. Na terceira linha, os "X" estão nas colunas 1, 3 e 5. As linhas 4 e 5 são idênticas, com "X" apenas nas colunas 1 e 5.} \\
            }
        }
    }
     \\ \\

    \theutterance \stepcounter{utterance}  
    & & & \multicolumn{2}{p{0.3\linewidth}}{
        \cellcolor[rgb]{0.9,0.9,0.9}{
            \makecell[{{p{\linewidth}}}]{
                \texttt{\tiny{[GM$|$GM]}}
                \texttt{continue} \\
            }
        }
    }
    & & \\ \\

    \theutterance \stepcounter{utterance}  
    & & \multicolumn{4}{p{0.6\linewidth}}{
        \cellcolor[rgb]{0.9,0.9,0.9}{
            \makecell[{{p{\linewidth}}}]{
                \texttt{\tiny{[GM$\rangle$P2]}}
                \texttt{Esta é a minha DESCRIÇÃODESCRIÇÃO:} \\
\texttt{A grelha é uma matriz 5x5 composta por "X" e quadrados vazios (▢). A estrutura forma um padrão simétrico:} \\
\texttt{{-} **1ª linha**: X ▢ ▢ ▢ X} \\
\texttt{{-} **2ª linha**: X X ▢ X X} \\
\texttt{{-} **3ª linha**: X ▢ X ▢ X} \\
\texttt{{-} **4ª linha**: X ▢ ▢ ▢ X} \\
\texttt{{-} **5ª linha**: X ▢ ▢ ▢ X} \\
\\ 
\texttt{As bordas laterais (1ª e 5ª colunas) são todas "X", exceto o centro da 3ª linha (posição 3,3). Há um "X" adicional no centro da 2ª linha (posição 2,2 e 2,4).} \\
\texttt{Agora faça uma pergunta para descobrir novos aspetos da minha grelha que possam ser diferentes da sua grelha. Comece com "PERGUNTA:" e não adicione mais nada.} \\
            }
        }
    }
    & \\ \\

    \theutterance \stepcounter{utterance}  
    & & & \multicolumn{4}{p{0.6\linewidth}}{
        \cellcolor[rgb]{0.9,0.9,0.9}{
            \makecell[{{p{\linewidth}}}]{
                \texttt{\tiny{[GM$\langle$P2]}}
                \texttt{PERGUNTA: Na sua grelha, a posição (3,3) — terceira linha, terceira coluna — está preenchida com "X" ou está vazia (▢)?} \\
            }
        }
    }
     \\ \\

    \theutterance \stepcounter{utterance}  
    & & & \multicolumn{2}{p{0.3\linewidth}}{
        \cellcolor[rgb]{0.9,0.9,0.9}{
            \makecell[{{p{\linewidth}}}]{
                \texttt{\tiny{[GM$|$GM]}}
                \texttt{continue} \\
            }
        }
    }
    & & \\ \\

    \theutterance \stepcounter{utterance}  
    & & \multicolumn{4}{p{0.6\linewidth}}{
        \cellcolor[rgb]{0.9,0.9,0.9}{
            \makecell[{{p{\linewidth}}}]{
                \texttt{\tiny{[P1$\langle$GM]}}
                \texttt{Esta é a minha DESCRIÇÃODESCRIÇÃO: A grelha é uma matriz 5x5 composta por "X" e quadrados vazios (▢). A estrutura forma um padrão simétrico com "X" nas bordas externas (primeira e última linha, primeira e última coluna), exceto nas posições centrais das linhas 2 a 5. Na segunda linha, há "X" nas colunas 1, 2, 4 e 5. Na terceira linha, os "X" estão nas colunas 1, 3 e 5. As linhas 4 e 5 são idênticas, com "X" apenas nas colunas 1 e 5.} \\
\texttt{PERGUNTA: Na sua grelha, a posição (3,3) — terceira linha, terceira coluna — está preenchida com "X" ou está vazia (▢)?Comece a sua resposta com "RESPOSTA:" e não acrescente mais nada.} \\
            }
        }
    }
    & \\ \\

    \theutterance \stepcounter{utterance}  
    & \multicolumn{4}{p{0.6\linewidth}}{
        \cellcolor[rgb]{0.9,0.9,0.9}{
            \makecell[{{p{\linewidth}}}]{
                \texttt{\tiny{[P1$\rangle$GM]}}
                \texttt{RESPOSTA: A posição (3,3) está preenchida com "X".} \\
            }
        }
    }
    & & \\ \\

    \theutterance \stepcounter{utterance}  
    & & & \multicolumn{2}{p{0.3\linewidth}}{
        \cellcolor[rgb]{0.9,0.9,0.9}{
            \makecell[{{p{\linewidth}}}]{
                \texttt{\tiny{[GM$|$GM]}}
                \texttt{continue} \\
            }
        }
    }
    & & \\ \\

    \theutterance \stepcounter{utterance}  
    & & \multicolumn{4}{p{0.6\linewidth}}{
        \cellcolor[rgb]{0.9,0.9,0.9}{
            \makecell[{{p{\linewidth}}}]{
                \texttt{\tiny{[P1$\langle$GM]}}
                \texttt{Agora faça uma pergunta para descobrir novos aspetos da minha grelha que possam ser diferentes da sua grelha. Comece com "PERGUNTA:" e não adicione mais nada.} \\
            }
        }
    }
    & \\ \\

    \theutterance \stepcounter{utterance}  
    & \multicolumn{4}{p{0.6\linewidth}}{
        \cellcolor[rgb]{0.9,0.9,0.9}{
            \makecell[{{p{\linewidth}}}]{
                \texttt{\tiny{[P1$\rangle$GM]}}
                \texttt{PERGUNTA: Na sua grelha, a posição (2,3) — segunda linha, terceira coluna — está preenchida com "X" ou está vazia (▢)?} \\
            }
        }
    }
    & & \\ \\

    \theutterance \stepcounter{utterance}  
    & & & \multicolumn{2}{p{0.3\linewidth}}{
        \cellcolor[rgb]{0.9,0.9,0.9}{
            \makecell[{{p{\linewidth}}}]{
                \texttt{\tiny{[GM$|$GM]}}
                \texttt{continue} \\
            }
        }
    }
    & & \\ \\

    \theutterance \stepcounter{utterance}  
    & & \multicolumn{4}{p{0.6\linewidth}}{
        \cellcolor[rgb]{0.9,0.9,0.9}{
            \makecell[{{p{\linewidth}}}]{
                \texttt{\tiny{[GM$\rangle$P2]}}
                \texttt{RESPOSTA: A posição (3,3) está preenchida com "X".} \\
\texttt{PERGUNTA: Na sua grelha, a posição (2,3) — segunda linha, terceira coluna — está preenchida com "X" ou está vazia (▢)?Comece a sua resposta com "RESPOSTA:" e não acrescente mais nada.} \\
            }
        }
    }
    & \\ \\

    \theutterance \stepcounter{utterance}  
    & & & \multicolumn{4}{p{0.6\linewidth}}{
        \cellcolor[rgb]{0.9,0.9,0.9}{
            \makecell[{{p{\linewidth}}}]{
                \texttt{\tiny{[GM$\langle$P2]}}
                \texttt{RESPOSTA: ▢} \\
            }
        }
    }
     \\ \\

    \theutterance \stepcounter{utterance}  
    & & & \multicolumn{2}{p{0.3\linewidth}}{
        \cellcolor[rgb]{0.9,0.9,0.9}{
            \makecell[{{p{\linewidth}}}]{
                \texttt{\tiny{[GM$|$GM]}}
                \texttt{continue} \\
            }
        }
    }
    & & \\ \\

    \theutterance \stepcounter{utterance}  
    & & \multicolumn{4}{p{0.6\linewidth}}{
        \cellcolor[rgb]{0.9,0.9,0.9}{
            \makecell[{{p{\linewidth}}}]{
                \texttt{\tiny{[GM$\rangle$P2]}}
                \texttt{Agora faça uma pergunta para descobrir novos aspetos da minha grelha que possam ser diferentes da sua grelha. Comece com "PERGUNTA:" e não adicione mais nada.} \\
            }
        }
    }
    & \\ \\

    \theutterance \stepcounter{utterance}  
    & & & \multicolumn{4}{p{0.6\linewidth}}{
        \cellcolor[rgb]{0.9,0.9,0.9}{
            \makecell[{{p{\linewidth}}}]{
                \texttt{\tiny{[GM$\langle$P2]}}
                \texttt{PERGUNTA: Na sua grelha, a posição (4,3) — quarta linha, terceira coluna — está preenchida com "X" ou está vazia (▢)?} \\
            }
        }
    }
     \\ \\

    \theutterance \stepcounter{utterance}  
    & & & \multicolumn{2}{p{0.3\linewidth}}{
        \cellcolor[rgb]{0.9,0.9,0.9}{
            \makecell[{{p{\linewidth}}}]{
                \texttt{\tiny{[GM$|$GM]}}
                \texttt{continue} \\
            }
        }
    }
    & & \\ \\

    \theutterance \stepcounter{utterance}  
    & & \multicolumn{4}{p{0.6\linewidth}}{
        \cellcolor[rgb]{0.9,0.9,0.9}{
            \makecell[{{p{\linewidth}}}]{
                \texttt{\tiny{[P1$\langle$GM]}}
                \texttt{RESPOSTA: ▢} \\
\texttt{PERGUNTA: Na sua grelha, a posição (4,3) — quarta linha, terceira coluna — está preenchida com "X" ou está vazia (▢)?Comece a sua resposta com "RESPOSTA:" e não acrescente mais nada.} \\
            }
        }
    }
    & \\ \\

    \theutterance \stepcounter{utterance}  
    & \multicolumn{4}{p{0.6\linewidth}}{
        \cellcolor[rgb]{0.9,0.9,0.9}{
            \makecell[{{p{\linewidth}}}]{
                \texttt{\tiny{[P1$\rangle$GM]}}
                \texttt{RESPOSTA: ▢} \\
            }
        }
    }
    & & \\ \\

    \theutterance \stepcounter{utterance}  
    & & & \multicolumn{2}{p{0.3\linewidth}}{
        \cellcolor[rgb]{0.9,0.9,0.9}{
            \makecell[{{p{\linewidth}}}]{
                \texttt{\tiny{[GM$|$GM]}}
                \texttt{continue} \\
            }
        }
    }
    & & \\ \\

    \theutterance \stepcounter{utterance}  
    & & \multicolumn{4}{p{0.6\linewidth}}{
        \cellcolor[rgb]{0.9,0.9,0.9}{
            \makecell[{{p{\linewidth}}}]{
                \texttt{\tiny{[P1$\langle$GM]}}
                \texttt{Agora faça uma pergunta para descobrir novos aspetos da minha grelha que possam ser diferentes da sua grelha. Comece com "PERGUNTA:" e não adicione mais nada.} \\
            }
        }
    }
    & \\ \\

    \theutterance \stepcounter{utterance}  
    & \multicolumn{4}{p{0.6\linewidth}}{
        \cellcolor[rgb]{0.9,0.9,0.9}{
            \makecell[{{p{\linewidth}}}]{
                \texttt{\tiny{[P1$\rangle$GM]}}
                \texttt{PERGUNTA: Na sua grelha, as posições (4,2) e (4,4) — quarta linha, segunda e quarta colunas — estão preenchidas com "X" ou vazias (▢)?} \\
            }
        }
    }
    & & \\ \\

    \theutterance \stepcounter{utterance}  
    & & & \multicolumn{2}{p{0.3\linewidth}}{
        \cellcolor[rgb]{0.9,0.9,0.9}{
            \makecell[{{p{\linewidth}}}]{
                \texttt{\tiny{[GM$|$GM]}}
                \texttt{continue} \\
            }
        }
    }
    & & \\ \\

    \theutterance \stepcounter{utterance}  
    & & \multicolumn{4}{p{0.6\linewidth}}{
        \cellcolor[rgb]{0.9,0.9,0.9}{
            \makecell[{{p{\linewidth}}}]{
                \texttt{\tiny{[GM$\rangle$P2]}}
                \texttt{RESPOSTA: ▢} \\
\texttt{PERGUNTA: Na sua grelha, as posições (4,2) e (4,4) — quarta linha, segunda e quarta colunas — estão preenchidas com "X" ou vazias (▢)?Comece a sua resposta com "RESPOSTA:" e não acrescente mais nada.} \\
            }
        }
    }
    & \\ \\

    \theutterance \stepcounter{utterance}  
    & & & \multicolumn{4}{p{0.6\linewidth}}{
        \cellcolor[rgb]{0.9,0.9,0.9}{
            \makecell[{{p{\linewidth}}}]{
                \texttt{\tiny{[GM$\langle$P2]}}
                \texttt{RESPOSTA: ▢ e ▢} \\
            }
        }
    }
     \\ \\

    \theutterance \stepcounter{utterance}  
    & & & \multicolumn{2}{p{0.3\linewidth}}{
        \cellcolor[rgb]{0.9,0.9,0.9}{
            \makecell[{{p{\linewidth}}}]{
                \texttt{\tiny{[GM$|$GM]}}
                \texttt{continue} \\
            }
        }
    }
    & & \\ \\

    \theutterance \stepcounter{utterance}  
    & & \multicolumn{4}{p{0.6\linewidth}}{
        \cellcolor[rgb]{0.9,0.9,0.9}{
            \makecell[{{p{\linewidth}}}]{
                \texttt{\tiny{[GM$\rangle$P2]}}
                \texttt{Agora faça uma pergunta para descobrir novos aspetos da minha grelha que possam ser diferentes da sua grelha. Comece com "PERGUNTA:" e não adicione mais nada.} \\
            }
        }
    }
    & \\ \\

    \theutterance \stepcounter{utterance}  
    & & & \multicolumn{4}{p{0.6\linewidth}}{
        \cellcolor[rgb]{0.9,0.9,0.9}{
            \makecell[{{p{\linewidth}}}]{
                \texttt{\tiny{[GM$\langle$P2]}}
                \texttt{PERGUNTA: Na sua grelha, a posição (5,2) — quinta linha, segunda coluna — está preenchida com "X" ou vazia (▢)?} \\
            }
        }
    }
     \\ \\

    \theutterance \stepcounter{utterance}  
    & & & \multicolumn{2}{p{0.3\linewidth}}{
        \cellcolor[rgb]{0.9,0.9,0.9}{
            \makecell[{{p{\linewidth}}}]{
                \texttt{\tiny{[GM$|$GM]}}
                \texttt{continue} \\
            }
        }
    }
    & & \\ \\

    \theutterance \stepcounter{utterance}  
    & & \multicolumn{4}{p{0.6\linewidth}}{
        \cellcolor[rgb]{0.9,0.9,0.9}{
            \makecell[{{p{\linewidth}}}]{
                \texttt{\tiny{[P1$\langle$GM]}}
                \texttt{RESPOSTA: ▢ e ▢} \\
\texttt{PERGUNTA: Na sua grelha, a posição (5,2) — quinta linha, segunda coluna — está preenchida com "X" ou vazia (▢)?Comece a sua resposta com "RESPOSTA:" e não acrescente mais nada.} \\
            }
        }
    }
    & \\ \\

    \theutterance \stepcounter{utterance}  
    & \multicolumn{4}{p{0.6\linewidth}}{
        \cellcolor[rgb]{0.9,0.9,0.9}{
            \makecell[{{p{\linewidth}}}]{
                \texttt{\tiny{[P1$\rangle$GM]}}
                \texttt{RESPOSTA: ▢} \\
            }
        }
    }
    & & \\ \\

    \theutterance \stepcounter{utterance}  
    & & & \multicolumn{2}{p{0.3\linewidth}}{
        \cellcolor[rgb]{0.9,0.9,0.9}{
            \makecell[{{p{\linewidth}}}]{
                \texttt{\tiny{[GM$|$GM]}}
                \texttt{continue} \\
            }
        }
    }
    & & \\ \\

    \theutterance \stepcounter{utterance}  
    & & \multicolumn{4}{p{0.6\linewidth}}{
        \cellcolor[rgb]{0.9,0.9,0.9}{
            \makecell[{{p{\linewidth}}}]{
                \texttt{\tiny{[P1$\langle$GM]}}
                \texttt{Agora faça uma pergunta para descobrir novos aspetos da minha grelha que possam ser diferentes da sua grelha. Comece com "PERGUNTA:" e não adicione mais nada.} \\
            }
        }
    }
    & \\ \\

    \theutterance \stepcounter{utterance}  
    & \multicolumn{4}{p{0.6\linewidth}}{
        \cellcolor[rgb]{0.9,0.9,0.9}{
            \makecell[{{p{\linewidth}}}]{
                \texttt{\tiny{[P1$\rangle$GM]}}
                \texttt{PERGUNTA: Na sua grelha, a posição (5,3) — quinta linha, terceira coluna — está preenchida com "X" ou vazia (▢)?} \\
            }
        }
    }
    & & \\ \\

    \theutterance \stepcounter{utterance}  
    & & & \multicolumn{2}{p{0.3\linewidth}}{
        \cellcolor[rgb]{0.9,0.9,0.9}{
            \makecell[{{p{\linewidth}}}]{
                \texttt{\tiny{[GM$|$GM]}}
                \texttt{continue} \\
            }
        }
    }
    & & \\ \\

    \theutterance \stepcounter{utterance}  
    & & \multicolumn{4}{p{0.6\linewidth}}{
        \cellcolor[rgb]{0.9,0.9,0.9}{
            \makecell[{{p{\linewidth}}}]{
                \texttt{\tiny{[GM$\rangle$P2]}}
                \texttt{RESPOSTA: ▢} \\
\texttt{PERGUNTA: Na sua grelha, a posição (5,3) — quinta linha, terceira coluna — está preenchida com "X" ou vazia (▢)?Comece a sua resposta com "RESPOSTA:" e não acrescente mais nada.} \\
            }
        }
    }
    & \\ \\

    \theutterance \stepcounter{utterance}  
    & & & \multicolumn{4}{p{0.6\linewidth}}{
        \cellcolor[rgb]{0.9,0.9,0.9}{
            \makecell[{{p{\linewidth}}}]{
                \texttt{\tiny{[GM$\langle$P2]}}
                \texttt{RESPOSTA: ▢} \\
\\ 
\texttt{*(Conclusão: As duas grelhas são **idênticas**! Todas as posições verificadas coincidem.)*} \\
            }
        }
    }
     \\ \\

    \theutterance \stepcounter{utterance}  
    & & & \multicolumn{2}{p{0.3\linewidth}}{
        \cellcolor[rgb]{0.9,0.9,0.9}{
            \makecell[{{p{\linewidth}}}]{
                \texttt{\tiny{[GM$|$GM]}}
                \texttt{continue} \\
            }
        }
    }
    & & \\ \\

    \theutterance \stepcounter{utterance}  
    & & \multicolumn{4}{p{0.6\linewidth}}{
        \cellcolor[rgb]{0.9,0.9,0.9}{
            \makecell[{{p{\linewidth}}}]{
                \texttt{\tiny{[GM$\rangle$P2]}}
                \texttt{Agora tome uma decisão. O que você acha: a sua grade e a outra grade descrita são a mesma?  Escreva "DECISÃO: mesma grade" se você acha que elas são a mesma grade ou "DECISÃO: grades diferentes" se você acha que elas são grades diferentes. Não adicione mais nada.} \\
            }
        }
    }
    & \\ \\

    \theutterance \stepcounter{utterance}  
    & & & \multicolumn{4}{p{0.6\linewidth}}{
        \cellcolor[rgb]{0.9,0.9,0.9}{
            \makecell[{{p{\linewidth}}}]{
                \texttt{\tiny{[GM$\langle$P2]}}
                \texttt{DECISÃO: mesma grade} \\
            }
        }
    }
     \\ \\

    \theutterance \stepcounter{utterance}  
    & & & \multicolumn{2}{p{0.3\linewidth}}{
        \cellcolor[rgb]{0.9,0.9,0.9}{
            \makecell[{{p{\linewidth}}}]{
                \texttt{\tiny{[GM$|$GM]}}
                \texttt{continue} \\
            }
        }
    }
    & & \\ \\

    \theutterance \stepcounter{utterance}  
    & & & \multicolumn{2}{p{0.3\linewidth}}{
        \cellcolor[rgb]{0.9,0.9,0.9}{
            \makecell[{{p{\linewidth}}}]{
                \texttt{\tiny{[GM$|$GM]}}
                \texttt{success} \\
            }
        }
    }
    & & \\ \\

    \theutterance \stepcounter{utterance}  
    & & \multicolumn{4}{p{0.6\linewidth}}{
        \cellcolor[rgb]{0.9,0.9,0.9}{
            \makecell[{{p{\linewidth}}}]{
                \texttt{\tiny{[P1$\langle$GM]}}
                \texttt{RESPOSTA: ▢} \\
\\ 
\texttt{*(Conclusão: As duas grelhas são **idênticas**! Todas as posições verificadas coincidem.)*} \\
\texttt{Agora tome uma decisão. O que você acha: a sua grade e a outra grade descrita são a mesma?  Escreva "DECISÃO: mesma grade" se você acha que elas são a mesma grade ou "DECISÃO: grades diferentes" se você acha que elas são grades diferentes. Não adicione mais nada.} \\
            }
        }
    }
    & \\ \\

    \theutterance \stepcounter{utterance}  
    & \multicolumn{4}{p{0.6\linewidth}}{
        \cellcolor[rgb]{0.9,0.9,0.9}{
            \makecell[{{p{\linewidth}}}]{
                \texttt{\tiny{[P1$\rangle$GM]}}
                \texttt{DECISÃO: mesma grade} \\
            }
        }
    }
    & & \\ \\

    \theutterance \stepcounter{utterance}  
    & & & \multicolumn{2}{p{0.3\linewidth}}{
        \cellcolor[rgb]{0.9,0.9,0.9}{
            \makecell[{{p{\linewidth}}}]{
                \texttt{\tiny{[GM$|$GM]}}
                \texttt{continue} \\
            }
        }
    }
    & & \\ \\

    \theutterance \stepcounter{utterance}  
    & & & \multicolumn{2}{p{0.3\linewidth}}{
        \cellcolor[rgb]{0.9,0.9,0.9}{
            \makecell[{{p{\linewidth}}}]{
                \texttt{\tiny{[GM$|$GM]}}
                \texttt{success} \\
            }
        }
    }
    & & \\ \\

\end{supertabular}
}

\end{document}
